% ══════════════════════════════════════════════════════════════════════════════
% Lecture 19: Shortest Paths — Bellman-Ford, Dijkstra, and Problem Solving
% ══════════════════════════════════════════════════════════════════════════════

% ══════════════════════════════════════════════════════════════════════════════
\section*{Kruskal's Algorithm (Recap)}
\addcontentsline{toc}{subsection}{Kruskal's Algorithm (Recap)}
% ══════════════════════════════════════════════════════════════════════════════

\subsection{Algorithm Description}

Kruskal's algorithm builds a Minimum Spanning Tree (MST) by starting from an empty forest $T$ and
greedily maintaining it:

\begin{itemize}[noitemsep]
  \item At each step, add the cheapest edge that does not create a cycle.
  \item Repeat until all vertices are connected.
\end{itemize}

\textbf{Claim:} $T$ is always a sub-forest of some MST (proof by induction on the number of
edges/components in $T$).

% ─────────────────────────────────────────────────────────────────────────────
\subsection{Implementation with Union-Find}

In each iteration we need to check whether the cheapest edge creates a cycle.
This is equivalent to checking whether its two endpoints belong to the same connected component
$\Rightarrow$ use the \textbf{disjoint-set (union-find)} data structure.

\begin{itemize}[noitemsep]
  \item Initially, each vertex is its own singleton set.
  \item When edge $(u, v)$ is added to $T$, perform \textsc{Union} of the two components.
  \item To test for a cycle, compare \textsc{Find-Set}$(u)$ and \textsc{Find-Set}$(v)$.
\end{itemize}

% ─────────────────────────────────────────────────────────────────────────────
\subsection{Complexity Analysis}

\begin{complexity}{Kruskal's Algorithm}{kruskal}
\begin{itemize}[noitemsep]
  \item Initialize $A$: $O(1)$
  \item First for-loop: $|V|$ \textsc{Make-Set} operations
  \item Sort $E$: $O(|E| \log |E|)$
  \item Second for-loop: $O(|E|)$ \textsc{Find-Set} and \textsc{Union} operations
  \item Total: $O\bigl((|V|+|E|)\,\alpha(|V|)\bigr) + O(|E|\log|E|) = O(|E|\log|E|) = O(|E|\log|V|)$
\end{itemize}
\textbf{Note:} If edges are already sorted, the time reduces to $O(|E|\,\alpha(|V|))$,
which is almost linear.
\end{complexity}

% ══════════════════════════════════════════════════════════════════════════════
\section*{Problem Solving — MST Applications}
\addcontentsline{toc}{subsection}{Problem Solving — MST Applications}
% ══════════════════════════════════════════════════════════════════════════════

\begin{remark}{Problem 1 — Minimum Feedback Edge Set}{feedbackedge}
A \textbf{feedback edge set} of a graph $G$ is a subset $F \subseteq E$ such that every cycle
in $G$ contains at least one edge in $F$. Equivalently, removing all edges in $F$ makes $G$ acyclic.

\textbf{Goal:} Compute the minimum-weight feedback edge set.

\textbf{Key insight:} $F = E \setminus T$, where $T$ is a maximum spanning tree of $G$.
Equivalently, take all edges \emph{not} in the MST: the MST spans the graph without cycles,
so the remaining edges form a minimal set that covers all cycles.
\end{remark}

\begin{remark}{Problem 2 — MST Cut and Cycle Properties}{mstprops}
Let $G = (V, E)$ be a connected graph with distinct edge weights.
\begin{enumerate}[noitemsep]
  \item \textbf{Cut property:} For any partition of $V$ into two subsets, the minimum-weight
        edge crossing the cut is in every MST.
  \item \textbf{Cycle property:} The maximum-weight edge in any cycle of $G$ is \emph{not}
        in the MST.
\end{enumerate}
\end{remark}

% ══════════════════════════════════════════════════════════════════════════════
\section*{The Shortest Path Problem}
\addcontentsline{toc}{subsection}{The Shortest Path Problem}
% ══════════════════════════════════════════════════════════════════════════════

% ─────────────────────────────────────────────────────────────────────────────
\subsection{Definition}

\begin{definition}{Shortest Path Problem}{shortestpath}
\textbf{Input:} A directed graph $G = (V, E)$ with edge weights $w: E \to \mathbb{R}$.

The \textbf{weight of a path} $p = \langle v_0, v_1, \ldots, v_k \rangle$ is:
\[
  w(p) = \sum_{i=1}^{k} w(v_{i-1}, v_i)
\]

The \textbf{shortest-path weight} from $u$ to $v$ is:
\[
  \delta(u,v) = \begin{cases}
    \min\{w(p) : u \leadsto v\} & \text{if a path from } u \text{ to } v \text{ exists} \\
    \infty & \text{otherwise}
  \end{cases}
\]

A \textbf{shortest path} from $u$ to $v$ is any path $p$ with $w(p) = \delta(u,v)$.
\end{definition}

% ─────────────────────────────────────────────────────────────────────────────
\subsection{Problem Variants}

\begin{itemize}[noitemsep]
  \item \textbf{Single-source:} Find shortest paths from a source $s$ to every other vertex.
  \item \textbf{Single-destination:} Find shortest paths to a given destination from all vertices.
        Can be solved by single-source after reversing all edge directions.
  \item \textbf{Single-pair:} Find the shortest path from $u$ to $v$.
        No algorithm known that improves over single-source in the worst case.
  \item \textbf{All-pairs:} Find shortest paths for every pair $(u,v)$.
        Can be solved by running single-source for each vertex; better algorithms also exist.
\end{itemize}

% ══════════════════════════════════════════════════════════════════════════════
\section*{Negative Weights and Applications}
\addcontentsline{toc}{subsection}{Negative Weights and Applications}
% ══════════════════════════════════════════════════════════════════════════════

\subsection{Negative-Weight Edges}

We allow negative edge weights in the graph.

\begin{remark}{Negative-Weight Cycles}{negcycle}
If a negative-weight cycle is reachable from the source, shortest paths are undefined: we can
traverse the cycle repeatedly to obtain paths of arbitrarily negative weight (length $-\infty$).

\begin{itemize}[noitemsep]
  \item Some algorithms (e.g.\ Dijkstra's) only work with non-negative weights.
  \item Bellman-Ford handles negative weights and can detect negative cycles.
\end{itemize}
\end{remark}

\subsection{Why Negative Weights?}

\textbf{Application:} Buying and selling currency (arbitrage detection).
Model currencies as vertices and log-exchange rates as edge weights.
A negative cycle in this graph corresponds to an arbitrage opportunity
(a sequence of trades that yields a net profit).

% ══════════════════════════════════════════════════════════════════════════════
\section*{Bellman-Ford Algorithm}
\addcontentsline{toc}{subsection}{Bellman-Ford Algorithm}
% ══════════════════════════════════════════════════════════════════════════════

% ─────────────────────────────────────────────────────────────────────────────
\subsection{Setup and the Relax Operation}

For each vertex $v \in V$, maintain:
\begin{itemize}[noitemsep]
  \item $\ell(v)$: current upper estimate of the shortest-path length from $s$ to $v$.
  \item $\pi(v)$: predecessor of $v$ on the current best path from $s$.
\end{itemize}

\textbf{Initialization:} $\ell(s) = 0$, $\ell(v) = \infty$ for $v \neq s$, $\pi(v) = \text{NIL}$.

The key subroutine is \textbf{relaxation}:

\begin{algorithm}[H]
\caption{Relax Operation}
\begin{algorithmic}[1]
\Procedure{Relax}{$u, v, w$}
  \If{$\ell(v) > \ell(u) + w(u,v)$}
    \State $\ell(v) \gets \ell(u) + w(u,v)$ \AlgComment{Improve estimate}
    \State $\pi(v) \gets u$                  \AlgComment{Update predecessor}
  \EndIf
\EndProcedure
\end{algorithmic}
\end{algorithm}

\textbf{Intuition:} Can we shorten the path to $v$ by routing through $u$?
If $\ell(u) + w(u,v) < \ell(v)$, then yes.

% ─────────────────────────────────────────────────────────────────────────────
\subsection{Algorithm}

\begin{algorithm}[H]
\caption{Bellman-Ford Algorithm}
\begin{algorithmic}[1]
\Procedure{Bellman-Ford}{$G, w, s$}
  \State \textsc{Init-Single-Source}$(G, s)$ \AlgComment{$\ell(s)\gets 0$; $\ell(v)\gets\infty$; $\pi(v)\gets\text{NIL}$}
  \For{$i = 1$ to $|V| - 1$}
    \For{each edge $(u,v) \in E$}
      \State \Call{Relax}{$u, v, w$}
    \EndFor
  \EndFor
  \For{each edge $(u,v) \in E$} \AlgComment{Check for negative cycles}
    \If{$\ell(v) > \ell(u) + w(u,v)$}
      \State \Return \textbf{false} \AlgComment{Negative cycle detected}
    \EndIf
  \EndFor
  \State \Return \textbf{true}
\EndProcedure
\end{algorithmic}
\end{algorithm}

\textbf{Idea:} After $i$ iterations, the algorithm has found all shortest paths using at most $i$
edges. Since any simple path has at most $|V|-1$ edges, $|V|-1$ iterations suffice.

% ─────────────────────────────────────────────────────────────────────────────
\subsection{Correctness}

\begin{remark}{Key note}{bfnote}
Bellman-Ford is only guaranteed to return correct shortest paths if \textbf{no negative cycles}
are reachable from the source. It can, however, be used to \emph{detect} such cycles.
\end{remark}

\begin{theorem}{Optimal Substructure}{optsubstruct}
If $\langle s, v_1, \ldots, v_k, v_{k+1} \rangle$ is a shortest path from $s$ to $v_{k+1}$,
then $\langle s, v_1, \ldots, v_k \rangle$ is a shortest path from $s$ to $v_k$.

\textbf{Proof:} Suppose toward contradiction that there exists a shorter path $p'$ from $s$
to $v_k$. Then $p' \to (v_k, v_{k+1})$ would be shorter than the assumed shortest path to
$v_{k+1}$, a contradiction.
\end{theorem}

\begin{theorem}{Bellman-Ford Correctness}{bfcorrect}
\textbf{Invariant:} After the $i$-th iteration, $\ell(v)$ is at most the length of the
shortest path from $s$ to $v$ using at most $i$ edges.

\textbf{Proof by induction:}
\begin{itemize}[noitemsep]
  \item \textit{Base case:} After 0 iterations, $\ell(s) = 0$ and $\ell(v) = \infty$ for
        $v \neq s$. Trivially correct.
  \item \textit{Inductive step:} Consider any shortest path from $s$ to $v_{k+1}$ using at
        most $i$ edges. By optimal substructure, the prefix to $v_k$ is the shortest path using
        at most $i-1$ edges. By the induction hypothesis, $\ell(v_k) \le w(p)$ after the
        previous iteration. After relaxing $(v_k, v_{k+1})$ in the $i$-th iteration,
        $\ell(v_{k+1}) \le \ell(v_k) + w(v_k, v_{k+1}) \le$ (length of shortest path using
        at most $i$ edges).
\end{itemize}

If there are no negative cycles reachable from $s$, then for any vertex $v$ there is a
shortest path using at most $n-1$ edges (a path with $\ge n$ edges must contain a cycle of
non-negative weight that can be removed). Hence Bellman-Ford returns the correct answer after
$n-1$ iterations.
\end{theorem}

% ─────────────────────────────────────────────────────────────────────────────
\subsection{Detecting Negative Cycles}

\begin{theorem}{Negative Cycle Detection}{negcycledetect}
There is a negative-weight cycle reachable from the source $s$ \textbf{if and only if} at
least one $\ell$-value changes when a further (the $n$-th) iteration of Bellman-Ford is
executed.

\textbf{Proof ($\Rightarrow$):} From the correctness proof, if there are no negative cycles,
$\ell$-values do not change in the $n$-th iteration.

\textbf{Proof ($\Leftarrow$):} Suppose no $\ell$-value changes in the $n$-th iteration. Then
for all $(u,v) \in E$: $\ell(u) + w(u,v) \ge \ell(v)$. For any cycle
$v_0 \to v_1 \to \cdots \to v_{t-1} \to v_0$:
\[
  \sum_{i=1}^{t} \ell(v_i) \;\le\; \sum_{i=1}^{t}\bigl(\ell(v_{i-1}) + w(v_{i-1},v_i)\bigr)
  \;=\; \sum_{i=1}^{t}\ell(v_{i-1}) + \sum_{i=1}^{t} w(v_{i-1},v_i)
\]
The two $\ell$-sums cancel (cyclic), giving $0 \le \sum_{i=1}^{t} w(v_{i-1},v_i)$, so the
cycle is non-negative. Hence no negative cycle is reachable.
\end{theorem}

% ─────────────────────────────────────────────────────────────────────────────
\subsection{Complexity Analysis}

\begin{complexity}{Bellman-Ford Algorithm}{bellmanford}
\begin{itemize}[noitemsep]
  \item \textsc{Init-Single-Source}: $\Theta(|V|)$
  \item Nested for-loops: $|V|-1$ iterations, each relaxing all $|E|$ edges: $\Theta(|E| \cdot |V|)$
  \item Final negative-cycle check: $\Theta(|E|)$
\end{itemize}
\textbf{Total:} $\Theta(|E| \cdot |V|)$
\end{complexity}

% ─────────────────────────────────────────────────────────────────────────────
\subsection{Final Comments on Bellman-Ford}

\begin{itemize}[noitemsep]
  \item Can detect negative cycles: run for $n$ iterations and look for cycles in the
        "shortest-path tree" — they correspond to negative cycles.
  \item Easy to implement in \textbf{distributed} settings: each vertex repeatedly queries
        its neighbours for the best path. This makes it useful for routing protocols and
        dynamic networks.
\end{itemize}

% ─────────────────────────────────────────────────────────────────────────────
\subsection{Problem Solving — Ski Resort Sanity Check}

\begin{remark}{Exam-style Problem}{skires}
Lindsey Vonn's assistant has mapped a ski resort as a directed weighted graph $G = (V,E)$
where vertices are stations and edge weights are altitude differences (positive for uphill,
negative for downhill).

\textbf{Sanity check:} Starting from any station $A$, the total altitude change when returning
to $A$ should always be 0 — i.e.\ no cycle should have a strictly negative or positive total
weight.

\textbf{Algorithm:} Run Bellman-Ford from any source vertex.
\begin{itemize}[noitemsep]
  \item If the algorithm reports a negative cycle, the map is inconsistent.
  \item Since the graph models altitude, a positive cycle would also be invalid; check both
        directions (or equivalently, also check for negative cycles in the reversed graph).
\end{itemize}
\textbf{Running time:} $O(|V| \cdot |E|)$.
\end{remark}

% ══════════════════════════════════════════════════════════════════════════════
\section*{Dijkstra's Algorithm}
\addcontentsline{toc}{subsection}{Dijkstra's Algorithm}
% ══════════════════════════════════════════════════════════════════════════════

% ─────────────────────────────────────────────────────────────────────────────
\subsection{Algorithm for Non-Negative Weights}

Dijkstra's algorithm is a faster alternative when \textbf{all edge weights are non-negative}.

\textbf{Key ideas:}
\begin{itemize}[noitemsep]
  \item Only works with non-negative weights.
  \item Greedy and faster than Bellman-Ford.
  \item Similar to Prim's algorithm for MST (essentially a weighted BFS).
\end{itemize}

\textbf{Approach:} Maintain a set $S$ of vertices whose shortest-path distances are finalised.
Greedily grow $S$: at each step, add the vertex $v \notin S$ that is \emph{closest} to $s$, i.e.\
\[
  \underset{v \notin S}{\arg\min}\; \min_{u \in S}\bigl(\ell(u) + w(u,v)\bigr)
\]

\begin{algorithm}[H]
\caption{Dijkstra's Algorithm}
\begin{algorithmic}[1]
\Procedure{Dijkstra}{$G, w, s$}
  \State \textsc{Init-Single-Source}$(G, s)$
  \State $S \gets \emptyset$
  \State $Q \gets V$ \AlgComment{Min-priority queue keyed by $\ell(v)$}
  \While{$Q \neq \emptyset$}
    \State $u \gets$ \Call{Extract-Min}{$Q$}
    \State $S \gets S \cup \{u\}$
    \For{each vertex $v \in \text{Adj}[u]$}
      \State \Call{Relax}{$u, v, w$}
      \If{$\ell(v)$ decreased}
        \State \Call{Decrease-Key}{$Q, v, \ell(v)$}
      \EndIf
    \EndFor
  \EndWhile
\EndProcedure
\end{algorithmic}
\end{algorithm}

% ─────────────────────────────────────────────────────────────────────────────
\subsection{Correctness}

\begin{theorem}{Dijkstra's Correctness}{dijkstracorrect}
\textbf{Loop invariant:} At the start of each iteration, for all $v \in S$,
$\ell(v) = \delta(s,v)$ (the distance estimate is exact).

\textbf{Proof sketch:} Let $u$ be the vertex being added to $S$. Suppose $\ell(u) > \delta(s,u)$.
Let $y$ be the first vertex on a true shortest path to $u$ that is not yet in $S$, and $x$ its
predecessor in $S$. Since $u$ was chosen over $y$, we have $\ell(u) \le \ell(y)$. But
$\ell(y) = \delta(s,y) \le \delta(s,u) \le \ell(u)$, so $\ell(u) = \delta(s,u)$,
contradicting the assumption.

\textbf{Why non-negative weights are required:} Once $u$ is added to $S$, its distance is
considered final. A negative-weight edge could later yield a shorter path, violating the invariant.
\end{theorem}

% ─────────────────────────────────────────────────────────────────────────────
\subsection{Implementation and Complexity}

Like Prim's, Dijkstra's performance depends on the priority queue implementation.

\begin{complexity}{Dijkstra's Algorithm}{dijkstra}
\begin{itemize}[noitemsep]
  \item With a \textbf{binary min-heap}: each operation takes $O(\log |V|)$, giving total
        $O(|E| \log |V|)$.
  \item With a \textbf{Fibonacci heap}: $O(|V| \log |V| + |E|)$.
\end{itemize}
\end{complexity}

% ─────────────────────────────────────────────────────────────────────────────
\subsection{Comparison: Bellman-Ford vs.\ Dijkstra}

\begin{center}
\begin{tabular}{lll}
\toprule
\textbf{Algorithm} & \textbf{Time complexity} & \textbf{Handles negative weights?} \\
\midrule
Bellman-Ford & $O(|V| \cdot |E|)$ & Yes (also detects negative cycles) \\
Dijkstra     & $O(|E| \log |V|)$  & No (requires non-negative weights) \\
\bottomrule
\end{tabular}
\end{center}

% ─────────────────────────────────────────────────────────────────────────────
\subsection{Problem Solving — Dijkstra Correctness Proof}

\begin{remark}{Exercise}{dijkstraexo}
Show that Dijkstra's algorithm is correct by proving the following loop invariant:

\medskip
\textit{``At the start of each iteration, for all $v \in S$, the distance $\ell(v)$ from $s$
to $v$ is equal to the shortest-path distance $\delta(s,v)$.''} \\

The proof proceeds by induction on $|S|$ as sketched in the correctness theorem above.
\end{remark}

% ─────────────────────────────────────────────────────────────────────────────
\subsection{Lecture Summary}

\begin{tcolorbox}[colback=lightblue, colframe=keyblue, fonttitle=\bfseries,
                  title={Key Points}]
\begin{itemize}[noitemsep]
  \item \textbf{Kruskal's Algorithm:} Greedy MST algorithm using union-find. Time:
        $O(|E|\log|V|)$; almost linear if edges are pre-sorted.
  \item \textbf{Shortest Path Problem:} Find minimum-weight paths from a source to all other
        vertices in a weighted directed graph.
  \item \textbf{Relaxation:} Core operation: improve $\ell(v)$ if routing through $u$ is
        cheaper.
  \item \textbf{Bellman-Ford:} Handles negative weights; detects negative cycles.
        Time: $\Theta(|V| \cdot |E|)$. Correct after $|V|-1$ iterations when no negative
        cycles exist.
  \item \textbf{Negative Cycle Detection:} A negative cycle is reachable iff some $\ell$-value
        decreases in the $n$-th iteration.
  \item \textbf{Dijkstra's Algorithm:} Greedy, faster alternative for non-negative weights.
        Uses a min-priority queue (analogous to Prim's MST). Time: $O(|E|\log|V|)$ with a
        binary heap.
  \item \textbf{Correctness of Dijkstra:} Maintains the invariant that each finalised vertex
        has its exact shortest-path distance. Fails with negative weights.
  \item \textbf{Choice:} Use Bellman-Ford when negative weights or cycle detection are needed;
        use Dijkstra for better performance on non-negative graphs.
\end{itemize}
\end{tcolorbox}